
D'après (1.11.3), on
peut écrire
@@M0@@
avec @@M1@@. On prend @@M2@@, ce qui est bien @@M3@@ si @@M4@@. On obtient alors :
@@M5@@

Pour @@M6@@, @@M7@@, donc d'après (1.11.2) :
@@M8@@

Par conséquent, la relation (1.11.4) se réduit à :
@@M9@@

Or @@M10@@ est de rang @@M11@@, donc @@M12@@ est la classe d'un faisceau inversible, et est donc inversible dans @@M13@@. Par suite, la relation (1.11.4) se réduit à la relation (1.11.2) pour @@M14@@.

Donc la relation (1.11.2) est vraie pour @@M15@@; pour @@M16@@, elle donne :
@@M17@@
Comme @@M18@@, on obtient @@M19@@. On montre ainsi par récurrence descendante sur @@M20@@ que les @@M21@@ sont tous nuls, en utilisant les relations (1.11.2) pour @@M22@@ variant de @@M23@@ à @@M24@@. Ceci montre que @@M25@@ sont indépendants sur @@M26@@, et achève la démonstration du théorème 1.1.

\textbf{Remarque 1.12.} Pour tout @@M27@@-Module localement libre de type fini @@M28@@, soit @@M29@@ la classe dans @@M30@@ de la @@M31@@-ième puissance symétrique de @@M32@@; posons @@M33@@~; les relations (1.11.3) peuvent encore s'écrire :
